\section{Methodology}

\subsection{Hypothesis Structure}

We formulate our research question as a structured hypothesis with testable predictions:

\textbf{Core Hypothesis (H-StaticFeedback-v1):} Under iterative LLM code repair on HumanEval/MBPP, integrating static analyzer feedback (pylint/mypy) alongside execution feedback improves pass@k compared to execution-only feedback, because static analysis provides mechanistic error explanations.

\textbf{Causal Mechanism:}
\begin{enumerate}
    \item Static analyzer produces warnings on LLM-generated code
    \item LLM receives mechanistic feedback (why, not just where)
    \item Targeted repairs lead to higher pass@k
\end{enumerate}

\textbf{Key Assumptions:}
\begin{itemize}
    \item A1: Static analyzers produce meaningful warnings on LLM code
    \item A2: LLMs can interpret static warnings into fixes
    \item A3: Static and execution feedback catch non-overlapping errors
    \item A4: HumanEval contains problems with static-detectable errors
    \item A5: Pylint severity filtering (disable conventions) is effective
\end{itemize}

\subsection{Gate-Based Validation}

We employ a hierarchical gate structure:

\begin{table}[h]
\centering
\caption{Gate structure for hypothesis validation.}
\label{tab:gates}
\begin{tabular}{@{}llll@{}}
\toprule
Gate & Type & Hypothesis & Criterion \\
\midrule
h-e1 & MUST\_WORK & Existence & $\geq$30\% of problems have actionable warnings \\
h-m1 & MUST\_WORK & Mechanism & LLM interprets warnings into fixes \\
h-m2 & SHOULD\_WORK & Mechanism & Static/execution non-overlapping \\
h-c1 & SHOULD\_WORK & Comparison & Improvement stratified by warnings \\
\bottomrule
\end{tabular}
\end{table}

MUST\_WORK gates block downstream experiments if failed. This design catches methodology issues early.

\subsection{Experimental Design}

\textbf{Variables:}
\begin{itemize}
    \item Independent: Feedback type (execution-only vs.\ static+execution)
    \item Dependent: pass@1, pass@5, pass@10
    \item Controlled: Base LLM, iteration budget (5), prompt format, static analyzer (pylint)
\end{itemize}

\textbf{Dataset:} HumanEval (164 problems)

\textbf{Static Analysis Configuration:}
\begin{verbatim}
pylint_config = {
    "disabled_categories": ["C", "R"],  # Convention, Refactoring
    "enabled_types": ["error", "warning"],
    "output_format": "json"
}
\end{verbatim}

\subsection{Existence Gate Protocol}

Before full pass@k evaluation, we validate the existence assumption: do static analyzers produce actionable warnings on benchmark code?

\textbf{Protocol:}
\begin{enumerate}
    \item Load all HumanEval problems
    \item Extract code (canonical solutions or LLM generations)
    \item Run pylint with configured filters
    \item Count problems with $\geq$1 actionable warning
    \item Gate passes if warning\_rate $\geq$ 30\%
\end{enumerate}

\textbf{Rationale:} If most solutions have zero warnings, static analysis provides no signal---the intervention is null. We set the 30\% threshold based on practical considerations: below this rate, fewer than 1 in 3 problems would receive static feedback, limiting the intervention's statistical power and practical utility. This threshold is conservative; lower rates (e.g., 20\%) might still enable meaningful experiments but would require larger sample sizes to detect effects.

\subsection{Metrics}

\begin{itemize}
    \item \textbf{Warning Rate:} Fraction of problems with $\geq$1 actionable warning
    \item \textbf{Warnings per Problem:} Average actionable warnings
    \item \textbf{Warning Distribution:} Breakdown by pylint code
\end{itemize}
